\documentclass[twocolumn,english,aps,prl,floatfix]{revtex4-1}
\setcounter{secnumdepth}{3}
\usepackage{color}
\usepackage{babel}
\usepackage{float}
\usepackage{amsmath}
\usepackage{amssymb}
\usepackage{graphicx}
\usepackage{esint}
\usepackage{MnSymbol}
%\usepackage[unicode=true,pdfusetitle,
% bookmarks=true,bookmarksnumbered=false,bookmarksopen=false,
% breaklinks=true,pdfborder={0 0 0},backref=false,colorlinks=true]
% {hyperref}
%\hypersetup{
% linkcolor=blue,citecolor=blue,urlcolor=blue}
%\usepackage{breakurl}

\makeatletter
%%%%%%%%%%%%%%%%%%%%%%%%%%%%%% User specified LaTeX commands.
%Format of numbering
\renewcommand{\theequation}{S\arabic{equation}}
\renewcommand{\thefigure}{S\arabic{figure}}

\makeatother

\newcommand{\w}{\omega}
\newcommand{\mimp}{m_{\rm imp}}

\newcommand{\todo}[1]{{\color{red} {\bf [ToDo: #1]}}}
\newcommand{\mv}[1]{{\color{green}{#1}}}

%%%%%%%%%%%%%%%%%%%%%%%%%%%%%%%%%%%%%%%%%%%%%%%%%%%%%%%%%%%%%%%%%%%%

\begin{document}

\title{Supplementary information for:\\
Destruction of long-range order in non-collinear two-dimensional antiferromagnets by random-bond disorder
%Destruction of long-range order in chiral two-dimensional antiferromagnets by random-bond disorder
}

\author{Santanu Dey}
\author{Matthias Vojta}
\affiliation{Institut f\"ur Theoretische Physik, Technische Universit\"at Dresden,
01062 Dresden, Germany}

\date{\today}

\maketitle
%%%%%%%%%%%%%%%%%%%%%%%%%%%%%%%%%%%%%%%%%%%%%%%%%%%%%%%%%%%%%%%%%%%%%%%

\section{Single bond defect}

Here we provide a few details concerning the physics of a single 
defect bond.

\subsection{Spin texture via linear response}

As explained in the main text, the spin texture induced by a bond 
defect can be understood as the response to a local dipolar 
transverse field which acts to rotate the two spins on the defect 
bond away from their clean-limit configuration. This is most 
efficiently calculated in a local frame where the clean-limit state 
is uniform. 

To this end we start from a coplanar state in the $x-y$ plane and 
rotate the spins such that all point along the $\hat{x}$ axis. The 
rotation involves the $y$ components of the spins, i.e., the 
susceptibility $\chi^\perp(\vec{q}) = \llangle S_{\vec{q}}^y; 
S^y_{-\vec{q}} \rrangle$. 
For a system with underlying SU(2) spin symmetry, this susceptibility
is governed by the in-plane Goldstone modes and generically behaves 
as $\chi^\perp(\vec{q})\propto 1/q^2$ for small momenta. For the 
triangular-lattice Heisenberg model in the classical limit the 
full expression for the transverse susceptibility is given by
\cite{wollny-thesis16}

\begin{align}
    \chi^\perp(\vec{q}) = S\frac{1}{\gamma(\vec{q})-1}
\end{align}
where the triangular lattice structure factor 

\begin{equation}\label{eq:gamma}
    \begin{aligned}
        \gamma(\vec{q})
        &=\frac{1}{6}\sum_{\vec{\delta}\in\text{NN}}
        e^{i\vec{q}\cdot\vec{\delta}} \\
        &=\frac{1}{3}\left(\cos q_x + 2\cos\frac{q_x}{2}
        \cos\frac{\sqrt{3}}{2}q_y\right)
    \end{aligned}
\end{equation}
encodes the geometry of the lattice.
For a bond defect at a position $\vec{r}'$ on a coupling along the 
direction $\vec{\delta}(\vec{r}')$, the transverse response at a 
separate site $\vec{r}$ is given by

\begin{align}
    \langle S_{\vec{r}}^\perp \rangle = \sum_{j=0}^1 h_{\vec{r}'}^\perp 
    \beta_j \llangle S_{\vec{r}}^y; 
    S_{r'+\beta_j\vec{\delta}(\vec{r}')/2}^y \rrangle 
\end{align}
where $\beta_j=(-1)^j$ is the dipolar form factor of the perturbation, 
and $h^\perp \propto \delta J/J$ measures the strength of the 
perturbation. The proportionality factor, which we call 
$\tilde{\alpha}$, depends on the microscopic details of the lattice 
and in itself is proportional to the stiffness of the ordered phase. 
Later, we extract the values of this non-universal prefactor for 
two different non-collinear antiferromagnets, the triangular-lattice 
Heisenberg model and the $J_1-J_4$ cubic-lattice Heisenberg model.

Once the transverse response is expressed in terms of the 
susceptibility $\chi^\perp(q)$,

\begin{align}
    \langle S_{\vec{r}}^\perp \rangle = \frac{\tilde{\alpha}}{J}
    \sum_{\vec{q}} \delta J(\vec{r}')
    \tilde{\beta}_{\vec{\delta}(\vec{r}')}(\vec{q})\chi^\perp(\vec{q})
    e^{i\vec{q}\cdot\left(\vec{r}-\vec{r}'\right)}
\end{align}
with the momentum dependence of the form factor given by 
$\tilde{\beta}_{\vec{\delta}(\vec{r}')}(\vec{q})=
e^{-i\vec{q}\cdot\vec{\delta}(\vec{r}')/2}-
e^{i\vec{q}\cdot\vec{\delta}(\vec{r}')/2}
=-2i\sin(\vec{q}\cdot\vec{\delta}(\vec{r}')/2)$, its dipolar nature 
becomes apparent. The long wave length response, 

\begin{align}\label{eq:tr_resp}
    \langle S_{\vec{r}}^\perp\rangle = \frac{\alpha S}{J}
    \sum_{\vec{q}}
    \delta J(\vec{r}')
    \frac{\vec{q}\cdot\vec{\delta}(\vec{r}')}{q^2}
    e^{i\vec{q}\cdot(\vec{r}-\vec{r}')}
\end{align}
and thereby the spin texture of the resulting state is 
determined by the defect induced dipole moment 
$\vec{d}(\vec{r}')=
(\delta J(\vec{r}')/J)\vec{\delta}
(\vec{r}')$ and the quadratic dispersion of the Goldstone 
modes. Once again the prefactor $\alpha$ (see also in main 
text) depends on the underlying lattice and following 
(\ref{eq:gamma}), for triangular lattice, we get 
$\alpha=4i\tilde{\alpha}$.

\subsection{Uniform moment}

In the presence of a single defect bond, the bulk spin moments do 
not exactly cancel each other, such that the resulting state has a 
non-vanishing uniform magnetization $\mimp$. For small $\delta J$ 
this is proportional to $\delta J$,
\begin{equation}
    \mimp = \gamma S \delta J + \mathcal{O}(S^0)\,.
\end{equation}
For the triangular-lattice Heisenberg model we numerically find 
$\gamma = $ \todo{add number}, see Fig.~\todo{add fig}. The 
proportionality can be expected to break down for larger $\delta J$, 
but still holds to very good accuracy for $\delta J=-J$ where we 
find $\mimp = 0.41 S + \mathcal{O}(S^0)$.
\todo{can you get the $1/S$ correction?}

A finite-size scaling analysis of $\mimp$ is shown in 
Fig.~\todo{add fig}.


%%%%%%%%%%%%%%%%%%%%%%%%%%%%%%%%%%%%%%%%%%%%%%%%%%%%%%%%%%%%%%%%%%%%%%%

\section{Destruction of long-range order}

As explained in the main text, an estimate of the transverse spin 
fluctuations induced by a collection of bond defects can be used to 
assess the destruction of long-range order.

In the linear response regime, we can consider the superposition of
individual dipolar responses (\ref{eq:tr_resp}) from a spatially
random distribution of bond defects. The average transverse response
is simply zero for a defect distribution with zero mean and therefore 
their collective effect is manifested through the squared transverse
fluctuation. Although the angular profile of the response due 
to each bond defect is dependent on the orientations
of the respective bonds, for the case of a perfectly random 
distribution of such defects, 
$\overline{\delta J(\vec{r})\delta J(\vec{r}')}=
(\delta J)^2\delta_{\vec{r},\vec{r}'}$, the contribution coming
from the separate bond or dipolar orientations decouple. From
(\ref{eq:tr_resp}) it is also easy to see that these contributions
differ by phase factors of constant angular rotations. Therefore by 
choosing appropriately rotated frames of axes for each, we can cast 
them in an identical form and sum up to find the expression for the 
total transverse spin fluctuation,


\begin{equation}\label{eq:sq_resp}
    \begin{aligned}
        \overline{\left\langle {S_0^{\perp}}^2\right\rangle } =& 
        \left(\frac{\alpha}{J}\right)^2
        \sum_{\vec{r},\vec{r}'}\sum_{\vec{q},\vec{q}'}
        \overline{\delta J(\vec{r})\delta J(\vec{r}')} \\
        &(\vec{q}\cdot\vec{\delta}(\vec{r}))
        (\vec{q}'\cdot\vec{\delta}(\vec{r}'))
        \chi^\perp(\vec{q})\chi^\perp(\vec{q}') 
        e^{i\left(\vec{q}\cdot\vec{r}+\vec{q}'\cdot\vec{r}'\right)}\\
        =& \left(\frac{\alpha\delta J}{J}\right)^2
        \sum_{\vec{r}}\sum_{\vec{q},\vec{q}'}
        (\vec{q}\cdot\vec{\delta}(\vec{r}))
        (\vec{q}'\cdot\vec{\delta}(\vec{r})) \\
        & \chi^\perp(\vec{q})\chi^\perp(\vec{q}') 
        e^{i\left(\vec{q}+\vec{q}'\right)\cdot\vec{r}}\\
        =& \left(\frac{\alpha\delta J}{J}\right)^2
        \frac{1}{N_B}\sum_{i}\sum_{\vec{q}}
        (\vec{q}\cdot\vec{\delta}_i)
        (\vec{q}\cdot\vec{\delta}_i)
        \chi^\perp(\vec{q})\chi^\perp(\vec{q}') \\
        =&\left(\frac{\alpha\delta J}{J}\right)^2 \int 
        \frac{d^d q}{(2\pi)^d}\left(q_x 
        \chi^\perp(\vec q) \right)^2
    \end{aligned}
\end{equation}
where $\chi^\perp(q)\propto S/q^2$ is the transverse susceptibility 
from the previous section and $N_B$ is the number of distinct bond 
orientation for a given lattice geometry. It is notworthy that the 
additional factor from the summation over the bond orientation is 
cancelled by the overall volume fraction for each such contribution. 
As discussed in the foremost part of this letter, if the susceptibility 
$\chi^\perp(\vec{q})$ is dominated by massless Goldstone modes, the 
integral (\ref{eq:sq_resp}) diverges below a lower critical dimension. 
This signal the destruction of the long-range order and precisely at 
the critical dimension emergence of exponentially growing domains. 
In the following we consider the criterion 
$ \overline{\left\langle {S^{\perp}}^2\right\rangle } \simeq S^2 $  
and provide a detailed supplement to our original discussion on 
various aspects of the destruction of longe-range order in 
bond-disordered non-collinear magnets.

\subsection{SU(2) symmetry and $T=0$}

For spontaneously broken SU(2) symmetry in a magnetic system 
the fundamental massless excitations have quadratic dispersions and 
therefore the transverse susceptibility 
$\chi^\perp(\vec{q})\propto 1/q^2$. The convergence criterion of the 
integral (\ref{eq:sq_resp}) makes the correlation length of 
the order parameter short-ranged,
\begin{equation}
    \begin{aligned}
        \left(\frac{\alpha S\delta J}{J}\right)^2
        \int_{1/\xi}^1q^{d-1}dq
        \left(\frac{1}{q}\right)^2
        & \simeq S^2
    \end{aligned}
\end{equation}
This implies a correlation length $\xi^{2-d} \propto (J/\delta J)^2$ 
for $d<2$, while for $d=2$, barring the non-universal stiffness 
dependent factor $\alpha$
\begin{equation}
    \xi \propto e^{(J/\delta J)^2}\,.
    \label{xi}
\end{equation}


\subsection{Finite $T$}

At finite $T$ and at $d=2$ the magnet would be thermally disordered,
but similar to the effect of bond disorder from the previous
segment is exhibits exponential correlation growth, 
$\xi \propto e^{1/T} $\cite{chakravarty87}. As temperature decreases,
the thermal and quenched bond disorder competes for precedence and for
a particular value of disorder strength a crossover temperature
emerges, $ T^* = \delta J/J$. If the temperature decreases any further
the disordered state becomes magnetic. \todo {This is apparently
contentious now, so we should change. In Chakravorty Halperin I 
see a factor of spin-stiffness in the exponent, should we try to be
precise and try get an appropriate cancellation of this with our
prefactor $\alpha$?
}

\subsection{Layered systems}

A realistic layered system displays a small inter-layer coupling 
which tends to stabilizes magnetic order at finite temperatures 
in the presence of SU(2) symmetry. We assume an (unfrustrated) 
inter-layer coupling of magnitude $J_\perp$. Using the ratio 
between the interplanar and intraplanar couplings,
$\kappa=J_\perp/J$ we can express the long-wavelength transverse 
susceptibility of the coupled layered system as

\begin{equation}
    \begin{aligned}
        \chi^\perp(\vec{q},q_\perp) &= 
        \frac{\lambda S}{q^2+\kappa^2q_\perp^2}
    \end{aligned}
\end{equation}
where as before we have a dimensionless lattice dependent 
factor $\lambda$. Substituting this expression in the integral 
relation (\ref{eq:sq_resp}) it becomes apparent that only the 
impurities among the in-plane bonds affect the longe-range order.
The stability criterion is determined by the coupling ratio 
$\kappa$ alone, 
\begin{equation}
    \begin{aligned}
    \left(\frac{\lambda\alpha S\delta J}{J}\right)^2
    \int \frac{d^2qdq_\perp}{(2\pi)^3}
    \frac{q_x}{q^2+\kappa^2q_\perp^2} &\leq S^2
    \end{aligned}
\end{equation}
with longe-range ordered phase stable if,
\begin{equation}
\begin{aligned}
    \frac{J_\perp}{J} = \kappa \geq 
    \frac{1}{2}\left(\frac{\lambda\alpha\delta J}{J}\right)^2
\end{aligned}
\end{equation}

\subsection{Broken SU(2)}

For a system with broken SU(2) symmetry, some or all of the 
Goldstone modes acquire a gap. For simplicity we assume a
particular form of the response function,
\begin{align}
    \chi^\perp(\vec{q}) = \frac{\lambda S}{q^2+\Delta^2}
\end{align}
with the dimensionless phenomenological factor $\lambda$
absorbing the particular details of the lattice. 

%%%%%%%%%%%%%%%%%%%%%%%%%%%%%%%%%%%%%%%%%%%%%%%%%%%%%%%%%%%%%%%%%%%%%%%

\section{Incommensurability in anisotropic Heisenberg model}

We consider a spatially anisotropic Heisenberg model on the triangular lattice, with couplings $J_{1,2,3}$ on the bonds $1,2,3$, with mean $J$. In the classical limit, the ground state is given by a single-$Q$ coplanar spiral whose propagation wavevector $\vec Q$ can be found by minimizing the Fourier-transformed coupling constant $J({\vec q})$. In the isotropic case, the system orders at $\vec Q_0 = \pm (4\pi/3,0)$.

The degree of anisotropy can be parameterized by $\vec{d} = \sum_{i=1}^3 (J_i/J) \hat{e}_i$ where $\hat{e}_i$ are the directed unit vectors along the nearest-neighbor bonds of the lattice. As explained in the main text, $\vec d$ takes the role of a local dipole density in the bond-disordered case.

For small deviations from isotropy, $|\vec d|\ll 1$, the relation between $\vec{Q}$ and $\vec d$ can be determined analytically \todo{work out, write down}.

%%%%%%%%%%%%%%%%%%%%%%%%%%%%%%%%%%%%%%%%%%%%%%%%%%%%%%%%%%%%%%%%%%%%%%%

\section{Numerical results}

%%%%%%%%%%%%%%%%%%%%%%%%%%%%%%%%%%%%%%%%%%%%%%%%%%%%%%%%%%%%%%%%%%%%%%%

\section{Bond defect in three dimensions}

We have checked that our scenario for bond disorder applies to three space dimensions as well. To this end, we have considered a single bond defect in \todo{describe model and results}.

\begin{thebibliography}{99}
    \bibitem{wollny-thesis16}
        %Fractional Moments and Singular Field Response
        A. Wollny,
        PhD Thesis. (2016).

    \bibitem{chakravarty87}
        %Low-temperature Behavior of Two-Dimensional Quantum
        %Antiferromagnets
        S. Chakravorty, B. I. Halperin, and D. R. Nelson
        Phys. Rev. Lett. \textbf{60}, 1057 (1988).

    \bibitem{chalker10}
        % Spin glass transition in geometrically frustrated antiferromagnets with weak disorder
        A. Andreanov, J. T. Chalker, T. E. Saunders, and D. Sherrington,
        Phys. Rev. B \textbf{81}, 014406 (2010).
\end{thebibliography}



\end{document} 
